\section*{Övningar}
\addcontentsline{toc}{section}{Övningar}

Lösningarna till uppgifterna som inte ber om kod står i \hyperref[ch:losningar]{appendix~A}.
Koden publiceras i kursrepot efter föreläsningen.

\exercise{Routing}{Resonemang}
\label{ex:3:1}
\begin{enumerate}[a)]
  \item Sker routingen i bussen eller i noden? Motivera med vad en fysisk ledning kan och inte kan.
  \item En nod tar emot en frame med korrekt checksumma, men \code{DST} är inte dess adress. Vad
    ska den göra?
  \item Varför ska den \emph{inte} skicka \code{Nack} i det läget?
\end{enumerate}

\exercise{Felmodellen}{Resonemang}
\label{ex:3:2}
För vart och ett av felen --- tappad byte, ändrad bit, fördröjning --- svara på:
\begin{enumerate}[a)]
  \item Blir framen strukturellt korrekt?
  \item Vilken mekanism upptäcker felet?
  \item Vad står parsern i för tillstånd efteråt?
\end{enumerate}
Ett av de tre felen ska inte orsaka något problem alls. Vilket, och vad betyder det om det ändå
gör det?

\exercise{Förlorad kvittens}{Avkodning}
\label{ex:3:3}
Nod A skickar en \code{StatusRequest} med \code{SEQ 0x0007} till nod B. B tar emot den, behandlar
den, och skickar \code{Ack}. Kvittensen tappas.
\begin{enumerate}[a)]
  \item Vad gör A härnäst, och efter vad?
  \item Vad ser B då?
  \item Vad ska B göra, och vad ska B inte göra?
  \item Vad hade hänt om B varit tyst i stället?
\end{enumerate}

\exercise{Timeout-värdet}{Resonemang}
\label{ex:3:4}
\begin{enumerate}[a)]
  \item Vad händer om timeout är för kort? Ge två konsekvenser.
  \item Vad händer om den är för lång?
  \item Vilken information skulle behövas för att välja värdet rätt, och varför har en nod
    sällan den?
\end{enumerate}

\exercise{Kedjan}{Resonemang}
\label{ex:3:5}
Räkna upp de fem mekanismerna i tillförlitlighetskedjan, i ordning. För var och en: vilket problem
löser den, och vilket nytt problem skapar den?
